% \iffalse meta-comment
%
% hit-report-cover.dtx 是开题报告的封面与扉页
%
% \fi
%
% \section{开题报告的封面与扉页}
%
% \changes{v3.2a}{2026/09/27}{研究生报告首页标签的拉开字距按字走，\hologo{pdfLaTeX} 下一个汉字是几个记号}
% \changes{v3.2a}{2026/09/14}{深圳博士中期从 \file{hithesisartplus.cls} 并入（与
%   \file{hithesisart.cls} 九成代码逐字相同），封面自己加载 \pkg{tikzpagenodes} 那三个包；
%   首页改用部件变体，基准是研究生那一页；威海本科的开题与中期改用哈尔滨那一套；
%   本部硕士封面添加「学位论文」或「实践成果」；正文照终稿排的那一档目录页用罗马页码；
%   字段表格里放宽 \cs{hfuzz}；\cs{makecover} 记下「前置已排」}
% \changes{v3.2a}{2026/09/26}{填写线与签字线改用 \cs{docxrun} 的下划线：原件的填写线是带下划线
%   的粗体空格，线照 Times New Roman 粗体的 post 表；先前数学模式的 \cs{underline} 跟着内容的
%   深度走，各行高低不一}
% \changes{v3.2a}{2026/09/26}{研究生与哈尔滨、威海本科报告的首页照原件的段落序列排，不再用
%   \cs{parbox} 加 \env{tabular} 摆：原件那一页是一串段落，字段是标签加一串带下划线的半角空格。
%   研究生照研究生院现行的模板（学院（部）、学科/专业学位类别，“（学位论文）”或“（实践成果）”
%   一行，落款研究生院制），各校区一样；本科照教务处的中期报告原件，校名字标嵌在段里}
%
% 深圳博士中期那份封面拿 \cs{AddToShipoutPicture} 往整页上贴图形，要 \pkg{eso-pic}
% 与 \pkg{tikzpagenodes}。别的组合下这三个包一个都不装。
%    \begin{macrocode}
%<*report-cover>
\bool_if:NF \g__hit_final_bool
  {
\hit@if@shenzhen@doctor@interim
  {
    \RequirePackage { tikzpagenodes }
    \RequirePackage { eso-pic }
    \RequirePackage { tikz }
  }

%    \end{macrocode}
%
% \emph{研究生报告的首页照研究生院现行的模板排。} 那一页在原件里是一串段落，不是表：
% 两个空段、校名、空段、“博士学位开题报告”、空段、“题\ 目：”、“（学位论文／实践成果）”、
% 六个字段、五个空段、“研究生院制”。字段一段一个，标签是敲的字与半角空格，后面跟一串
% 加下划线的半角空格当填写线，学院那一行 26 个、学科 25 个、别的 29 个，所以填写线的起止
% 参差不齐，照抄。取值放在填写线正中。学号那个标签原件的 run 上有 0.6 磅字距，
% 每个字与空格各多两份（没有字符网格，平衡单字节与双字节字符开着）。
% 各段的格式在 \file{hit-layout.dtx}（\option{report-grad-} 那几个样式与 \option{cover}
% 那一节的页面设置），标签在 \file{hit-terms.dtx}。说明页原件接在首页后面，这里不排。
%
% 威海硕士开题照 \issue[219] 附的模板，没有“（学位论文／实践成果）”那一行：两个类型词条
% 都是空的，那一段就只排空段，文档里填了 \option{report-type} 也不印。
%    \begin{macrocode}
\tl_new:N \l__hit_report_grad_label_tl
\cs_new_protected:Npn \__hit_report_grad_blank:n #1
  { \docxlike_format_par:nn {#1} { \mbox { } } }
\cs_new_protected:Npn \__hit_report_grad_label:nn #1#2
  {
    \tl_set:Ne \l__hit_report_grad_label_tl {#2}
    \hit@engine@group@chars:N \l__hit_report_grad_label_tl
    \dim_compare:nNnTF {#1} = { 0pt }
      { \l__hit_report_grad_label_tl }
      {
        \tl_map_inline:Nn \l__hit_report_grad_label_tl
          { ##1 \skip_horizontal:n { (#1) * 2 } }
      }
  }
\cs_new_protected:Npn \__hit_report_grad_field:nnnnn #1#2#3#4#5
  {
    \docxlike_format_par:nnn { report-grad-field } {#1}
      {
        \__hit_report_grad_label:nn {#2} {#3}
        \docxrun [ underline-style = single , font-style = bold ]
          { \hbox_to_wd:nn { 0.5em * #4 } { \hfil #5 \hfil } }
      }
  }
\cs_new_protected:Npn \__hit_report_grad_type:
  {
    \tl_clear:N \l__hit_report_grad_label_tl
    \bool_lazy_and:nnF
      { \tl_if_blank_p:e { \hit@term@cover@report@type@dissertation } }
      { \tl_if_blank_p:e { \hit@term@cover@report@type@practice } }
      { \__hit_report_grad_type_label: }
    \tl_if_blank:VF \l__hit_report_grad_label_tl { （ \l__hit_report_grad_label_tl ） }
  }
\cs_new_protected:Npn \__hit_report_grad_type_label:
  {
    \cs_if_exist:NT \hit@info@report@type@zh
      { \tl_set:Ne \l__hit_report_grad_label_tl { \hit@info@report@type@zh } }
    \tl_if_blank:VT \l__hit_report_grad_label_tl
      {
        \tl_set:Ne \l__hit_report_grad_label_tl
          {
            \bool_if:NTF \g__hit_practice_bool
              { \hit@term@cover@report@type@practice }
              { \hit@term@cover@report@type@dissertation }
          }
      }
  }
\cs_new_protected:Npn \hit@report@title@page@graduate
  {
    \hit@layout@push@part:n { cover }
    \group_begin:
      \dim_zero:N \parindent
      \dim_zero:N \parskip
      \__hit_report_grad_blank:n { report-grad-blank-head }
      \__hit_report_grad_blank:n { report-grad-blank-head }
      \docxlike_format_par:nn { report-grad-university }
        {
          \hit@term@cover@university@zh
          \hit@if@option@T { weihai } { （ \hit@term@cover@campus ） }
        }
      \__hit_report_grad_blank:n { report-grad-blank-head }
      \docxlike_format_par:nn { report-grad-name }
        {
          \hit@term@cover@degree@level@zh \hit@term@cover@document@type
          \hit@stage@name \hit@term@cover@stage@document@type
        }
      \__hit_report_grad_blank:n { report-grad-blank }
      \docxlike_format_par:nn { report-grad-title }
        {
          \hit@term@cover@title@prefix \hit@term@cover@title@separator@zh
          \hit@info@title@cover@zh
        }
      \docxlike_format_par:nn { report-grad-type } { \__hit_report_grad_type: }
      \__hit_report_grad_field:nnnnn { paragraph = { spacing = { before = 12bp } } }
        { 0pt } { \hit@term@cover@affiliation } { 26 } { \hit@info@affiliation@zh }
      \__hit_report_grad_field:nnnnn { } { 0pt } { \hit@term@cover@major } { 25 }
        { \hit@info@major@zh }
      \__hit_report_grad_field:nnnnn { } { 0pt } { \hit@term@cover@supervisor } { 29 }
        { \hit@info@supervisor@zh }
      \__hit_report_grad_field:nnnnn { } { 0pt } { \hit@term@cover@student } { 29 }
        { \hit@info@author@zh }
      \__hit_report_grad_field:nnnnn { } { 0.6bp } { \hit@term@cover@student@id } { 29 }
        { \hit@info@student@id }
      \__hit_report_grad_field:nnnnn { } { 0pt } { \hit@term@cover@report@date } { 29 }
        { \hit@info@date@zh }
      \prg_replicate:nn { 5 } { \__hit_report_grad_blank:n { report-grad-blank-foot } }
      \docxlike_format_par:nn { report-grad-mark } { \hit@term@cover@bottom@mark@graduate }
    \group_end:
    \hit@layout@pop@part:
  }
%    \end{macrocode}
%
% 旧版首页（\cs{parbox} 加 \env{tabular} 排的）留下的修订记录：
% \changes{v3.0.22}{2022/05/16}{深圳校区中期模板封面范例中为“哈尔滨工业大学”。\issue[126]}
% \changes{v3.1c}{2023/04/16}{修改深圳校区硕士开题封面中的字号与字距。\issue[141]}
% \changes{v3.1c}{2023/04/16}{修改深圳校区硕士中期封面中的字号与字距。\issue[166]}
% \changes{v3.1d}{2025/03/03}{修改威海校区硕士开题封面中的字号与字距。\issue[219]}
% \changes{v3.0.03}{2020/07/03}{fix author name bug}
% \changes{v3.0.22}{2022/05/16}{深圳校区中期模板封面范例中没有“哈工大（深圳）制”字样。\issue[126]}
% \changes{v3.1c}{2023/04/16}{深圳校区硕士开题封面范例中有“深圳校区教务部~制”}
%    \begin{macrocode}
%    \end{macrocode}
%
% \emph{哈尔滨、威海本科报告的首页照教务处的原件排}（《本科毕业论文（设计）中期报告》，开题那份
% 同一套、只差阶段）：一个 12 磅空段、两个 18 磅空段、嵌着校名字标的一段、18 磅空段、文种、两个
% 空段、“题\ 目：”、四个空段、五个字段、两个空段、落款。字段是标签、一个半角空格、21 个
% 加下划线的半角空格；专业名长过填写线时填写线跟着放长（\issue[199]）。格式在
% \file{hit-layout.dtx} 的 \option{report-bach-} 那几个样式。字标照轮廓画，框照原件
% 167.8 $\times$ 31.25\,磅，见 \file{src/engine/hit-outline.dtx}。
%    \begin{macrocode}
\dim_new:N \l__hit_report_bach_line_dim
\cs_new_protected:Npn \__hit_report_bach_field:nn #1#2
  {
    \docxlike_format_par:nn { report-bach-field }
      {
        #1 \hit@sp
        \docxrun [ underline-style = single , font-style = bold ]
          { \hbox_to_wd:nn { \l__hit_report_bach_line_dim } { \hfil #2 \hfil } }
      }
  }
\cs_new_protected:Npn \hit@report@title@page@bachelor
  {
    \hit@layout@push@part:n { cover }
    \group_begin:
      \dim_zero:N \parindent
      \dim_zero:N \parskip
      \__hit_report_grad_blank:n { report-bach-blank-small }
      \__hit_report_grad_blank:n { report-bach-blank-head }
      \__hit_report_grad_blank:n { report-bach-blank-head }
      \docxlike_format_par:nn { report-bach-blank-head }
        { \hit@wordmark:nn { 167.8bp } { 31.25bp } }
      \__hit_report_grad_blank:n { report-bach-blank-head }
      \docxlike_format_par:nn { report-bach-name }
        {
          \hit@term@cover@document@type@level \hit@term@cover@document@type
          \hit@stage@name \hit@term@cover@stage@document@type
        }
      \__hit_report_grad_blank:n { report-bach-blank }
      \__hit_report_grad_blank:n { report-bach-blank }
      \docxlike_format_par:nn { report-bach-title }
        {
          \hit@term@cover@title@prefix \hit@term@cover@title@separator@zh
          \hit@info@title@cover@zh
        }
      \__hit_report_grad_blank:n { report-bach-blank }
      \prg_replicate:nn { 3 } { \__hit_report_grad_blank:n { report-bach-blank-wide } }
      \docxlike_format_apply_font:nN { report-bach-field } \l_tmpa_tl
      \hbox_set:Nn \l_tmpa_box
        {
          \l_tmpa_tl
          \dim_gset:Nn \g_tmpa_dim { 10.5em }
          \bfseries \hit@info@major@zh
        }
      \dim_set:Nn \l__hit_report_bach_line_dim
        { \dim_max:nn { \g_tmpa_dim } { \box_wd:N \l_tmpa_box } }
      \__hit_report_bach_field:nn { \hit@term@cover@major@bachelor } { \hit@info@major@zh }
      \__hit_report_bach_field:nn { \hit@term@cover@student@bachelor } { \hit@info@author@zh }
      \__hit_report_bach_field:nn { \hit@term@cover@student@id@bachelor } { \hit@info@student@id }
      \__hit_report_bach_field:nn { \hit@term@cover@supervisor@bachelor } { \hit@info@supervisor@zh }
      \__hit_report_bach_field:nn { \hit@term@cover@report@date@bachelor } { \hit@info@date@zh }
      \__hit_report_grad_blank:n { report-bach-blank }
      \__hit_report_grad_blank:n { report-bach-blank }
      \docxlike_format_par:nn { report-bach-mark } { \hit@cover@bottommark }
    \group_end:
    \hit@layout@pop@part:
  }
%    \end{macrocode}
%
% 旧版本科首页（\cs{parbox} 加 \env{tabular} 排的）留下的修订记录：
% \changes{v3.1d}{2025/03/03}{根据 \issue[213] 更改哈尔滨本科开题和中期封面，多了“本科”}
% \changes{v3.1d}{2025/03/03}{本科专业可能太长，自适应一下。\issue[199]}
%    \begin{macrocode}
\cs_new_protected:Npn \hit@report@back@page@bachelor
  {
    \thispagestyle { empty }
    \noindent \parbox [t] [6.5cm] [t] { \textwidth } { \hit@term@cover@teacher@comment@bachelor }
    \noindent \parbox [b] [6cm] [t] { \textwidth }
      {
        \hit@term@cover@teacher@signature@bachelor \docxrun [ underline-style = single ] { \makebox [3cm] { } }
        \hfill \hit@term@cover@check@date@bachelor \docxrun [ underline-style = single ] { \makebox [3cm] { } }
      }
  }

%    \end{macrocode}
%
% \changes{v3.0.22}{2022/05/16}{深圳校区中期模板目录中一级标题为黑体 \cs{heiti}，而不是加粗 \cs{bfseries}。\issue[126]}
% \changes{v3.1c}{2023/04/16}{修改错误的判断语句。\issue[141]}
% \changes{v3.1c}{2023/04/16}{深圳校区硕士开题目录有标题。\issue[166]}
% \changes{v3.1c}{2023/04/16}{深圳校区硕士开题目录修改样式}
% \changes{v3.1c}{2023/04/16}{深圳校区硕士中期目录修改样式。\issue[166]}
% \changes{v3.1d}{2025/03/03}{威海校区硕士开题目录修改样式。\issue[219]}
% 深圳校区中期模板目录中一级标题为黑体 \cs{heiti}。\par
% 深圳校区硕士开题目录有标题。
% 深圳校区硕士中期目录也有标题。
% 用 \pkg{tocloft} 宏包实现深圳校区硕士开题和中期的目录，逐步进行迁移。
% \begin{macro}{\makecover}
%
% \changes{v3.1c}{2023/04/16}{新增对空白页的控制。\issue[138]}
% \changes{v3.1c}{2023/04/16}{修改深圳硕士中期的页脚。\issue[166]}
% 封面分派。本科与研究生两套封面；\option{print} 打开时补一张空白页供双面打印；
% \option{toc} 打开时排目录。深圳硕士的页眉另有一套。
%    \begin{macrocode}

\cs_new_protected:Npn \hit@report@shenzhen@doctor@title@page@graduate
  {
~
 \hit@fontsize@x{xiaosi}{1}\textcolor[rgb]{1,1,1}{\songti{\hit@term@cover@hi}}%
 \vspace*{10mm}~
 \begin{center}~
 \songti\erhao\textbf{\hit@term@cover@university@zh\hit@if@option@T{shenzhen}{（\hit@term@cover@campus）}}
 \end{center}~
 \begin{center}~
 \songti\erhao\textbf{\hit@term@cover@check@report@title}
 \end{center}~
 \vspace{66bp}~
 \parbox[b][3cm][t]{\textwidth}{~
 \begin{center}\songti\sihao
 \cs_set:Npn \arraystretch {1.67}~
 \begin{tabular}{l@{\ \ }c}~
 \textbf{\hit@term@cover@student}~&~\textbf{\hit@info@author@zh}\\~
 \hline
 \textbf{\hit@term@cover@enroll@date}~&~\textbf{\hit@info@enroll@date@zh}\\~
 \hline
 \textbf{\hit@term@cover@title@label}~&~\textbf{\hit@info@title@line@one@zh}\\~
 \hline
 &\textbf{\hit@info@title@line@two@zh}\\~
 \hline
 \textbf{\hit@term@cover@affiliation@major}~&~\textbf{\hit@info@major@zh}\\~
 \hline
 \textbf{\hit@term@cover@report@date}~&~\textbf{\hit@info@date@zh}\\~
 \hline
 \textbf{\hit@term@cover@supervisor}~&~\textbf{\hit@info@supervisor@zh}\\~
 \hline
 \end{tabular}\cs_set:Npn \arraystretch {1}
 \end{center}~
 }~
  }

\cs_new_protected:Npn \hit@report@shenzhen@doctor@frame@picture
  {
%
~\begin{tikzpicture}[remember~picture,overlay]~
 \draw ([shift={(-5mm,0mm)}]current~page~text~area.south~west)~rectangle~([shift={(5mm,2mm)}]current~page~text~area.north~east);~
 \end{tikzpicture}%
  }
%    \end{macrocode}
%
% 深圳博士中期的封面自成一版：一张表格排完全部字段，正文四周画一个边框，
% 封面之后跟一页填写说明。
%    \begin{macrocode}
\cs_set_protected:Npn \hit@report@shenzhen@doctor@makecover
  {
    \hit@write@pdf@metadata
    \begin { titlepage }
      \hit@report@shenzhen@doctor@title@page@graduate
      \clearpage
      \hit@term@cover@note@interim@shenzhen@doctor
      \hit@structure@switch:nnn { table-of-contents }
        { \bool_if_p:N \g__hit_report_toc_auto_bool } { \clearpage \tableofcontents \clearpage }
    \end { titlepage }
    \AddToShipoutPicture { \hit@report@shenzhen@doctor@frame@picture }
    \pagestyle { hit@shenzhen@doctor@headings@main }
  }
%    \end{macrocode}
%
% \emph{目录那一页的页码。} 正文照终稿排的那一档与终稿一样分两段数：目录一段，
% 正文另起一段从 1 数（在 \file{src/flow/hit-report-matter.dtx} 换回来）。目录那一节
% 中期原件写着大写罗马（\texttt{w:fmt="upperRoman"}），开题两份（中、英）什么都没写，
% 就是阿拉伯数字，Windows Word 导出件的目录页印的正是“I”与“1”。首页不排页码，所以这一句
% 发在首页之后、目录之前。别的档整篇一个制式，不动。
%
% 那个开关一并发，而且要\emph{全局}发：页脚那一头按它分前置与正文两档装饰
% （见 \file{src/flow/hit-final-pagestyle.dtx} 的 \cs{hit@final@footer}），
% 报告从不切它，不发就一直按正文算，目录页会印成 \texttt{-I-}。原件目录页不带短横，
% 正文才是 \texttt{-1-}。
%
% \emph{局部赋值没用。} 页脚是页面送出时才求值的，那会儿 \env{titlepage} 这个
% 分组早关了，开关被还原成真，目录页照旧印带短横的。
%
% \emph{目录的字号行距。} 正文不照 Word 排的那几档，目录一直是小四单倍：先前是首页顶上那行
% 白字（\cs{hit@fontsize@x}\{xiaosi\}\{1\}）顺带设的，与目录同在 \env{titlepage} 这一组里。首页改照
% 原件排以后没有那行了，这里明写一遍，目录照旧。
%
% \emph{首页是个部件。} 基准那一页给研究生，本科与深圳本科各是一个变体：
% 本科那一份在本文件里，深圳本科那一份在
% \file{src/pages/hit-report-cover-shenzhen-bachelor.dtx}。取哪一份由档位定，
% 见 \file{src/engine/hit-key-engine.dtx} 的 \cs{hit@hook@new:nn}。
%    \begin{macrocode}
\hit@hook@new:nn { report-cover-page } { \hit@report@title@page@graduate }
\hit@hook@new:nn { report-cover-page-bachelor } { \hit@report@title@page@bachelor }
\cs_set_protected:Npn \hit@report@default@makecover
  {
    \hit@write@pdf@metadata
    \begin { titlepage }
      \hit@hook@use:n { report-cover-page }
      \clearpage
      \bool_if:NT \g__hit_print_bool
        { \thispagestyle { empty } \mbox { } \clearpage }
      \bool_if:NT \g__hit_final_body_bool
        {
          \legacy_if_gset_false:n { @mainmatter }
          \pagenumbering
            {
              \bool_lazy_and:nnTF
                { \bool_if_p:N \g__hit_shenzhen_bool } { \bool_if_p:N \g__hit_proposal_bool }
                { arabic } { Roman }
            }
          \pagestyle
            { \bool_if:NTF \g__hit_harbin_report_bool { hit@empty } { hit@headings@front } }
        }
      \bool_if:NF \g__hit_final_body_bool { \hit@fontsize@x { xiaosi } { 1 } }
      \hit@structure@switch:nnn { table-of-contents }
        { \bool_if_p:N \g__hit_report_toc_auto_bool }
        { \tableofcontents \hit@report@toc@tail: \clearpage }
    \end { titlepage }
    \bool_lazy_and:nnT
      { \bool_if_p:N \g__hit_shenzhen_bool } { \bool_if_p:N \g__hit_master_bool }
      { \pagestyle { hit@shenzhen@headings@main } }
  }
%    \end{macrocode}
%
% 排完置位。\file{hit-structure.dtx} 的兜底在文末发现没调过 \cs{hitfrontmatter}
% 就补调一次，而老写法的文档用的是 \cs{makecover}，不置位就会被补排一遍——
% 拿 v3.1e 的示例报告实测 3 页变 5 页，多出来的正是封面与目录。
%    \begin{macrocode}
\cs_new_protected:Npn \hit@report@makecover@mark:n #1
  {
    \bool_gset_true:N \g__hit_matter_front_bool
    \use:c {#1}
  }
\hit@if@shenzhen@doctor@interim@TF
  { \cs_set:Npn \makecover
      { \hit@report@makecover@mark:n { hit@report@shenzhen@doctor@makecover } } }
  { \cs_set:Npn \makecover
      { \hit@report@makecover@mark:n { hit@report@default@makecover } } }
%    \end{macrocode}
% \end{macro}
% \begin{macro}{\makebackcover}
%    \begin{macrocode}
\cs_set_protected:Npn \makebackcover
  { \clearpage \hit@report@back@page@bachelor }
  }
%</report-cover>
%    \end{macrocode}
% \end{macro}
%
% \Finale
